
\documentclass[11pt]{article}

\usepackage{amsmath}

\begin{document}

\section{Algebraic Theories}

In the Awodey-Bauer notes, the theory of groups is given by equations
(pp. 10 and 14)
\begin{subequations}
\begin{align}
   m(x, m(y, z)) & = m(m(x, y), z)
   \label{law:associative}
   \\
   m(x, e()) & = x
   \label{law:right-identity}
   \\
   m(e(), x) & = x
   \label{law:left-identity}
   \\
   m(x, i(x)) & = e()
   \label{law:right-inverse}
   \\
   m(i(x), x) & = e()
   \label{law:left-inverse}
\end{align}
\end{subequations}
where $m$ is a binary operation (the group operation, usually called
multiplication or division, according to taste), $i$ is a unary operation
(the inversion operation, $i(x)$ is the inverse of $x$), and $e$ is a nullary
operation, also called a constant, ($e()$ is the identity element of the
group).  There is no universe specified: all of variables range over the same
set or type.  If we call that $G$, then
\begin{align*}
   m & : G \times G \to G
   \\
   i & : G \to G
   \\
   e & : 1 \to G
\end{align*}
but since there is only one type the axioms do not need to mention it.
Also there is no quantification, at least not explicitly.  One can say
there is an implicit for-all for each variable, but not explicit.

The terms of the theory are the obvious.
\begin{itemize}
\item Each variable is a term.
\item Each function evaluated at any terms is a term.
\end{itemize}
Then the axioms are the equations that given to hold between terms.
The only rules for proof are variable renaming, terms for function evaluation,
and substution of equals for equals.

If we had fewer or more axioms (equations), then we would have different
theories.  For example,
\begin{itemize}
\item The theory with one binary function $m$ and one equation
    \eqref{law:associative} is the theory of \emph{semigroups}.
\item The theory with one binary function $m$, one nullary function $e$,
    and three equations \eqref{law:associative}, \eqref{law:right-identity},
    and \eqref{law:left-identity} is the theory of \emph{monoids},
    also called \emph{semigroups with identities}.
\item The theory with the same functions and equations as the theory
    of groups plus the additional equation
\begin{equation} \label{law:commutative}
    m(x, y) = m(y, x)
\end{equation}
    is the theory of \emph{commutative groups},
    also called \emph{abelian groups}.
\item The theory with the same functions and equations as the theory
    of groups plus the additional equation
\begin{equation} \label{law:boolean}
    m(x, x) = e
\end{equation}
    is the theory of \emph{boolean groups}.
\end{itemize}

\end{document}
